\documentclass[../../main.tex]{subfiles}
\begin{document}


{\bf [解]}
数$k$（$k=1,2,\dots,6$）の出る確率を$p_k$とおくと全体の確率の和は$1$だから
\begin{align}
 \sum_{k=1}^6p_k=1 \label{eq:1}
\end{align}
である．

\paragraph{(1)}

サイコロを$2$回ふったとき同じ目が出る確率$P$は
\begin{align}
 P=\displaystyle\sum_{k=1}^6p_k^2
\end{align}
である．
コーシー・シュワルツの不等式より
\begin{align}
 & \Bigl(\sum_{k=1}^6p_k\Bigr)^2\le6\sum_{k=1}^6p_k^2 \\
\therefore\
 & P\ge\dfrac{1}{6}\Bigl(\sum_{k=1}^6p_k\Bigr)^2=\dfrac{1}{6} \quad(\because\text{\cref{eq:1}}).
\end{align}
であるから題意の不等式は示された．
等号成立は，コーシー・シュワルツの等号成立条件より全ての$k$について$p_k=\dfrac{1}{6}$のときである．

\paragraph{(2)}

1回目に奇数，2回目に偶数の目が出る確率$Q$は
\begin{align}
 Q=(p_1+p_3+p_5)(p_2+p_4+p_6)
\end{align}
である．
$p_k\ge0$と\cref{eq:1}より，相加相乗平均の不等式を用いて
\begin{align}
  &\sqrt{Q}\le\dfrac{(p_1+p_3+p_5)+(p_2+p_4+p_6)}{2}=\dfrac{1}{2}\sum_{k=1}^6p_k = \dfrac{1}{2} \\
\therefore\ 
  & Q\le\dfrac{1}{4}\quad(\because Q>0) \label{eq:2}
\end{align}
であるから，題意の左側の不等号は示された．

次に右側の不等号を考える．
$t=p_1+p_3+p_5$（$0\le t\le1$）とおくと\cref{eq:1}より$1-t=p_2+p_4+p_6$である．
コーシー・シュワルツの不等式より
\begin{align}
 \begin{cases}
  3(p_1^2+p_3^2+p_5^2) \ge (p_1+p_3+p_5)^2 = t^2 \\
  3(p_2^2+p_4^2+p_6^2) \ge (p_2+p_4+p_6)^2 = (1-t)^2
 \end{cases}
\end{align}
である．辺々加えて(1)の$P$の定義を用いて
\begin{align}
  & 3P \ge t^2 + (1-t)^2 \\
\therefore
  & P\ge\dfrac{t^2+(1-t)^2}{3}
\end{align}
だから
\begin{align}
 \dfrac{1}{2}-\dfrac{3}{2}P\ 
  & \le\ \dfrac{1}{2}-\dfrac{1}{2}\bigl\{t^2+(1-t)^2\bigr\} \\
  & =\dfrac{1}{2}\bigl[1-(2t^2-2t+1)\bigr] \\
  & =t(1-t)
  & =Q \\
\therefore
  \dfrac{1}{2}-\dfrac{3}{2}P \le Q \label{eq:3}
\end{align}
である．よって右側の不等号が示された．
以上\cref{eq:2,eq:3}より
\begin{align}
\dfrac{1}{2}-\dfrac{3}{2}P\ \le\ Q\ \le\ \dfrac{1}{4}.
\end{align}
であって題意は示された．


\newpage

\end{document}
